% Lösungen Einheit 1 von 5 – Proportionale Funktion (zusammenbau v0.1)
\einheitenkopf{Einheit 1 von 5 · Proportionale Funktion}

\erg{9}{a) 3 (zum Beispiel $3 : 1 = 3$) \quad b) y-Werte: $0$, $3$, $6$; Gerade durch $(0|0)$, $(1|3)$, $(2|6)$ \quad c) $y = 6$; Rechnung: $2 \cdot 3 = 6$ \quad d) y-Werte: $0$, $3$, $6$; Gerade durch $(0|0)$, $(2|3)$, $(4|6)$ \quad e) y-Werte: $0$, $1$, $2$; Gerade durch $(0|0)$, $(3|1)$, $(6|2)$ \quad f) y-Werte: $0$, $-3$, $-6$; Gerade durch $(0|0)$, $(1|-3)$, $(2|-6)$; die Gerade fällt \quad g) Ja: alle Quotienten $y : x = 7$, also $y = 7x$ \quad h) Ben hat die Differenz statt des Quotienten genommen; richtig: $10 : 2 = 5$, also $y = 5x$. \quad i) $y = 12x$; Gerade durch $(0|0)$ und $(5|60)$; der Faktor 12 heißt: je Minute kommen 12 Liter dazu.}
\erg{10}{a) 15 € (1 m kostet 2{,}50 €; $2{,}5 \cdot 6 = 15$ passt)}
\erg{11}{a) x und y sind vertauscht; richtig ist $(2|6)$. \quad b) Nein: Die Gerade einer proportionalen Funktion geht immer durch den Ursprung. \quad c) $y = 1{,}5x$ (denn $3 : 2 = 1{,}5$) \quad d) Ja: $1{,}7 \cdot 35 = 59{,}50$ €, das ist weniger als 60 €.}
